\chapter{Fazit}

In diesem Versuch wird wieder einmal mehr die Bedeutung der Regel- und Messtechnik veranschaulicht. 
Mit ein paar wenigen Bauteilen lassen sich also Herzschläge von Menschen messen. 
Nötig dafür sind lediglich Operationsverstärker und ein paar Filter, um ein besseres Signal zu erhalten. 
Unsere Messergebnisse sind in Ordnung. 
Nur die Messung des Tiefpasses und der Bandsperre lieferte uns keine ausreichenden Ergebnisse. 
So haben wir unseren eigenen Hochpass gebaut und diesen vermessen. 
Deutlich erkennbar, war der Unterschied in der Frequenz der Peaks. 
Nach Digitaler Analyse war dann auch der klassische Herzschlag als solcher zu sehen. 
Der Versuch ist weitgehend an Elektrotechnik orientiert, man verbringt die meiste Zeit damit,
die einzelnen Bauteile auf einem Schaltbrett herumzustecken. 
Es war dennoch sehr interessant und gibt grundlegendes Verständnis der Regel- und Messtechnik.\\
\begin{figure}[H]
	\centering
	\includegraphics[width=0.6\textwidth]{Abb/herzschlag.jpg}
	\caption{Gute Messung eines Herzschlags}
\end{figure}
%Fänd ich noch schön als vergleichsbild zu unseren Herzschlägen.
%https://de.depositphotos.com/114828470/stock-illustration-heartbeat-icon-vector-illustration.html
